% AUTO-GENERATED by docs/diagrams/render.py from data/processed/processed_traces.jsonl.
% ProcessBench trace processbench:gsm8k:gsm8k-147 (generator: Qwen2-1.5B-Instruct). Do not edit by hand.
\begin{tikzpicture}[font=\sffamily\scriptsize, >={Stealth[length=2mm]},
  tracestep/.style={rectangle, rounded corners=1.5pt, draw=black!60, align=left,
               text width=12.5cm, inner sep=2mm},
  okstep/.style={tracestep, fill=green!7}, badstep/.style={tracestep, fill=red!10, draw=red!70!black, line width=0.8pt},
  afterstep/.style={tracestep, fill=gray!8, draw=black!35, dashed}]
\node[tracestep, fill=blue!5] (q) at (0,0) {\textbf{Question} (gsm8k tier): Jame will turn 27 in 5 years.  In 8 years his cousin will be 5 years younger than twice his age.  How many years separate the age of the two now?};
\node[okstep, below=2.5mm of q] (s1) {\textbf{$s_{1}$}: If Jame will turn 27 in 5 years, then currently he is 27 - 5 = 22 years old.};
\draw[->] (q.south) -- (s1.north);
\node[badstep, below=2.5mm of s1] (s2) {\textbf{$s_{2}$} \textcolor{red!70!black}{\textbf{first error: $e=2$, $\rho=e/n=0.50$}}: In 8 years, his cousin will be 2 * 22 - 8 = 46 - 8 = 38 years old.};
\draw[->] (s1.south) -- (s2.north);
\node[afterstep, below=2.5mm of s2] (s3) {\textbf{$s_{3}$} \textit{(after the error)}: To find out how many years separate the age of the two now, we need to subtract their current ages from each other: 38 - 22 = 16.};
\draw[->] (s2.south) -- (s3.north);
\node[afterstep, below=2.5mm of s3] (s4) {\textbf{$s_{4}$} \textit{(after the error)}: The answer is: \$\textbackslash{}boxed\{16\}\$};
\draw[->] (s3.south) -- (s4.north);
\node[tracestep, fill=black!6, below=2.5mm of s4] (a) {\textbf{Outcome}: final answer incorrect ($y=1$); $n=4$ steps, 2 step(s) remain after the first error.};
\draw[->] (s4.south) -- (a.north);
\end{tikzpicture}
